% GENERATED FILE. Edit canonical lessons, then run npm run generate:books.
% canonical-source-hash: fnv1a64:c6f5f6032be9c3ea
% canonical-lessons: AR-C410-abaha, AR-C410-awjaba, AR-C410-khalafa, AR-C410-iltazama, AR-C410-afa-exempt

\chapter{Allowed or required}
\label{ch:ar-allowed-or-required}
\hlchaptermodality{410}

\begin{chapteropening}
\textbf{By the end of this chapter:} \emph{I can say that something is permitted, obligatory or exempt, and that someone broke or kept a rule.}

Complete the last lesson of chapter 410: I can say that something is permitted, obligatory or exempt, and that someone broke or kept a rule.
\end{chapteropening}

\section[\texorpdfstring{\ar{أباح} (abāḥa)}{أباح (abāḥa)}]{\ar{أباح} (abāḥa) — to permit}

\label{lesson:AR-C410-abaha}



\begin{quote}
Before the new one: say the Arabic for to stipulate in writing, then the Arabic for to legislate.
\end{quote}

\subsection*{You'll want to know: \ar{أباح}}
\textbf{\ar{أباح}} — \emph{abāḥa} — \textquotedblleft{}to permit\textquotedblright{}.

\begin{quote}
\ar{أباح} \ar{القانون} \ar{العمل} \ar{يوم} \ar{السبت}.
\end{quote}

\begin{quote}
\emph{abāḥa al-qānūn al-ʿamal yawm as-sabt.}
\end{quote}

\begin{quote}
The law permitted work on Saturdays.
\end{quote}

Present: \ar{يبيح}, \emph{yubīḥ}.

\begin{cousinweb}
\textbf{\ar{أباح}} (\emph{abāḥa}) is built on the native Arabic root \emph{b-w-ḥ}, \textquotedblleft{}to become known, to disclose\textquotedblright{} (Lane). It is Form IV, to lay a thing open to people: to permit it.
\end{cousinweb}

\begin{grammarlens}[title={What you've built}]
One more word for saying what happens if.
\end{grammarlens}

\subsection*{Your turn}
\begin{itemize}
\raggedright
  \item \emph{Say it:} \emph{abāḥa}
  \item \emph{Say it:} \emph{abāḥa}, once more
  \item \emph{Say it:} \emph{abāḥa al-qānūn al-ʿamal yawm as-sabt.} — The law permitted work on Saturdays.
  \item \emph{Recall:} say the Arabic for to stipulate in writing, then the Arabic for to legislate, then say \emph{abāḥa} again
\end{itemize}

\begin{tcolorbox}[breakable,colback=teal!4,colframe=teal!35!black,title={Before you move on}]
What does \ar{أباح} mean? (\textquotedblleft{}To permit\textquotedblright{}.) Say it once more.
\end{tcolorbox}
\section[\texorpdfstring{\ar{أوجب} (awjaba)}{أوجب (awjaba)}]{\ar{أوجب} (awjaba) — to make obligatory}

\label{lesson:AR-C410-awjaba}



\begin{quote}
Before the new one: say the Arabic for to legislate, then the Arabic for to permit.
\end{quote}

\subsection*{You'll want to know: \ar{أوجب}}
\textbf{\ar{أوجب}} — \emph{awjaba} — \textquotedblleft{}to make obligatory\textquotedblright{}.

\begin{quote}
\ar{أوجب} \ar{القانون} \ar{التعليم} \ar{للأطفال}.
\end{quote}

\begin{quote}
\emph{awjaba al-qānūn at-taʿlīm li-l-aṭfāl.}
\end{quote}

\begin{quote}
The law made schooling obligatory for children.
\end{quote}

Present: \ar{يوجب}, \emph{yūjib}.

\begin{cousinweb}
\textbf{\ar{أوجب}} (\emph{awjaba}) is built on the native Arabic root \emph{w-j-b}, \textquotedblleft{}to be due, binding\textquotedblright{} (Lane). It is Form IV, to make binding: to make obligatory. \ar{واجب}, a duty, is from the same root.
\end{cousinweb}

\begin{grammarlens}[title={What you've built}]
One more word for saying what happens if.
\end{grammarlens}

\subsection*{Your turn}
\begin{itemize}
\raggedright
  \item \emph{Say it:} \emph{awjaba}
  \item \emph{Say it:} \emph{awjaba}, once more
  \item \emph{Say it:} \emph{awjaba al-qānūn at-taʿlīm li-l-aṭfāl.} — The law made schooling obligatory for children.
  \item \emph{Recall:} say the Arabic for to legislate, then the Arabic for to permit, then say \emph{awjaba} again
\end{itemize}

\begin{tcolorbox}[breakable,colback=teal!4,colframe=teal!35!black,title={Before you move on}]
What does \ar{أوجب} mean? (\textquotedblleft{}To make obligatory\textquotedblright{}.) Say it once more.
\end{tcolorbox}
\section[\texorpdfstring{\ar{خالف} (khālafa)}{خالف (khālafa)}]{\ar{خالف} (khālafa) — to break a rule}

\label{lesson:AR-C410-khalafa}



\begin{quote}
Before the new one: say the Arabic for to permit, then the Arabic for to make obligatory.
\end{quote}

\subsection*{You'll want to know: \ar{خالف}}
\textbf{\ar{خالف}} — \emph{khālafa} — \textquotedblleft{}to break a rule\textquotedblright{}.

\begin{quote}
\ar{خالف} \ar{الرجل} \ar{قواعد} \ar{المكان}.
\end{quote}

\begin{quote}
\emph{khālafa ar-rajul qawāʿid al-makān.}
\end{quote}

\begin{quote}
The man broke the rules of the place.
\end{quote}

Present: \ar{يخالف}, \emph{yukhālif}.

Also to differ: \ar{خالفه} \ar{الرأي}, disagreed with him.

\begin{cousinweb}
\textbf{\ar{خالف}} (\emph{khālafa}) is built on the native Arabic root \emph{kh-l-f}, \textquotedblleft{}to come after, behind\textquotedblright{} (Lane). It is Form III, to go against the one behind or before you: to break a rule, to differ.
\end{cousinweb}

\begin{grammarlens}[title={What you've built}]
One more word for saying what happens if.
\end{grammarlens}

\subsection*{Your turn}
\begin{itemize}
\raggedright
  \item \emph{Say it:} \emph{khālafa}
  \item \emph{Say it:} \emph{khālafa}, once more
  \item \emph{Say it:} \emph{khālafa ar-rajul qawāʿid al-makān.} — The man broke the rules of the place.
  \item \emph{Recall:} say the Arabic for to permit, then the Arabic for to make obligatory, then say \emph{khālafa} again
\end{itemize}

\begin{tcolorbox}[breakable,colback=teal!4,colframe=teal!35!black,title={Before you move on}]
What does \ar{خالف} mean? (\textquotedblleft{}To break a rule\textquotedblright{}.) Say it once more.
\end{tcolorbox}
\section[\texorpdfstring{\ar{التزم} (iltazama)}{التزم (iltazama)}]{\ar{التزم} (iltazama) — to comply}

\label{lesson:AR-C410-iltazama}



\begin{quote}
Before the new one: say the Arabic for to make obligatory, then the Arabic for to break a rule.
\end{quote}

\subsection*{You'll want to know: \ar{التزم}}
\textbf{\ar{التزم}} — \emph{iltazama} — \textquotedblleft{}to comply\textquotedblright{}.

\begin{quote}
\ar{التزم} \ar{الجميع} \ar{بالقواعد}.
\end{quote}

\begin{quote}
\emph{iltazama al-jamīʿ bi-l-qawāʿid.}
\end{quote}

\begin{quote}
Everyone complied with the rules.
\end{quote}

Present: \ar{يلتزم}, \emph{yaltazim}.

It takes \ar{بـ}: \ar{التزم} \ar{بـ} ….

\begin{cousinweb}
\textbf{\ar{التزم}} (\emph{iltazama}) is built on the native Arabic root \emph{l-z-m}, \textquotedblleft{}to cling, to stick to\textquotedblright{} (Lane). It is Form VIII, to bind oneself to a thing: to comply, to commit.
\end{cousinweb}

\begin{grammarlens}[title={What you've built}]
One more word for saying what happens if.
\end{grammarlens}

\subsection*{Your turn}
\begin{itemize}
\raggedright
  \item \emph{Say it:} \emph{iltazama}
  \item \emph{Say it:} \emph{iltazama}, once more
  \item \emph{Say it:} \emph{iltazama al-jamīʿ bi-l-qawāʿid.} — Everyone complied with the rules.
  \item \emph{Recall:} say the Arabic for to make obligatory, then the Arabic for to break a rule, then say \emph{iltazama} again
\end{itemize}

\begin{tcolorbox}[breakable,colback=teal!4,colframe=teal!35!black,title={Before you move on}]
What does \ar{التزم} mean? (\textquotedblleft{}To comply\textquotedblright{}.) Say it once more.
\end{tcolorbox}
\section[\texorpdfstring{\ar{أعفى} (aʿfā)}{أعفى (aʿfā)}]{\ar{أعفى} (aʿfā) — to exempt}

\label{lesson:AR-C410-afa-exempt}



\begin{quote}
Before the new one: say the Arabic for to break a rule, then the Arabic for to comply.
\end{quote}

\subsection*{You'll want to know: \ar{أعفى}}
\textbf{\ar{أعفى}} — \emph{aʿfā} — \textquotedblleft{}to exempt\textquotedblright{}.

\begin{quote}
\ar{أعفى} \ar{المدير} \ar{الموظف} \ar{من} \ar{العمل} \ar{يوم} \ar{العيد}.
\end{quote}

\begin{quote}
\emph{aʿfā al-mudīr al-muwaẓẓaf min al-ʿamal yawm al-ʿīd.}
\end{quote}

\begin{quote}
The manager exempted the employee from work on the feast day.
\end{quote}

Present: \ar{يعفي}, \emph{yuʿfī}.

It takes \ar{من}: \ar{أعفاه} \ar{من} ….

\begin{cousinweb}
\textbf{\ar{أعفى}} (\emph{aʿfā}) is built on the native Arabic root \emph{ʿ-f-w}, \textquotedblleft{}to efface; to pardon\textquotedblright{} (Lane). It is Form IV, to release someone from a thing: to exempt.
\end{cousinweb}

\begin{grammarlens}[title={What you've built}]
One more word for saying what happens if.
\end{grammarlens}

\subsection*{Your turn}
\begin{itemize}
\raggedright
  \item \emph{Say it:} \emph{aʿfā}
  \item \emph{Say it:} \emph{aʿfā}, once more
  \item \emph{Say it:} \emph{aʿfā al-mudīr al-muwaẓẓaf min al-ʿamal yawm al-ʿīd.} — The manager exempted the employee from work on the feast day.
  \item \emph{Recall:} say the Arabic for to break a rule, then the Arabic for to comply, then say \emph{aʿfā} again
\end{itemize}

\begin{tcolorbox}[breakable,colback=teal!4,colframe=teal!35!black,title={Before you move on}]
What does \ar{أعفى} mean? (\textquotedblleft{}To exempt\textquotedblright{}.) Say it once more.
\end{tcolorbox}
